% =====================================================================
%  Chương 2 — CƠ SỞ LÝ THUYẾT  (bản viết ngắn gọn, bám sát repo)
%  Gói cần trong main.tex: amsmath, amssymb, bbm, booktabs, float,
%                           hyperref, tikz + arrows.meta/positioning
% =====================================================================
\providecommand{\cmark}{\ensuremath{\checkmark}}
\providecommand{\xmark}{\ensuremath{\times}}

\chapter{II. Cơ sở lý thuyết}
\label{chap:coso_lythuyet}

Chương này tóm tắt các nền tảng lý thuyết được sử dụng trực tiếp trong
hệ thống: bài toán ABSA và các bài toán con, kiến trúc Transformer,
mô hình tiền huấn luyện BERT/T5 kèm kỹ thuật giải mã và chuẩn hóa đầu
ra, cơ chế Local Context Focus (LCF) và kỹ thuật Token Merging (ToMe).
Mỗi mục tập trung vào các khái niệm và công thức cốt lõi; chi tiết
hiện thực được trình bày ở Chương~\ref{chap:phuongphap}.

% =====================================================================
\section{2.1.\; Phân tích cảm xúc theo khía cạnh (ABSA)}
\label{sec:absa}

\subsection{2.1.1.\; Định nghĩa và các bài toán con}
\label{subsec:absa_subtasks}

\textbf{ABSA} (Aspect-Based Sentiment Analysis) phân tích cảm xúc
\emph{theo từng khía cạnh} thay vì toàn bộ câu. Câu
\textit{``The room was beautiful but the staff was rude''} mang cảm
xúc \textbf{Positive} với \textit{room} và \textbf{Negative} với
\textit{staff} -- phân tích cảm xúc mức câu không biểu diễn được sự
mâu thuẫn này \cite{pontiki2016semeval}.

ABSA bao gồm một nhóm bài toán con (Bảng~\ref{tab:absa_subtasks}).
Khóa luận tập trung vào hai bài toán: \textbf{ATE} (trích xuất cụm
từ khía cạnh) và \textbf{APC} (phân loại cảm xúc + nhóm khía cạnh);
các bài toán end-to-end phức tạp hơn (ASTE, ASQP) chỉ được khảo sát
ở mức tổng quan.

\begin{table}[H]
\centering
\caption{Các bài toán con chính của ABSA}
\label{tab:absa_subtasks}
\small
\begin{tabular}{@{}llll@{}}
\toprule
\textbf{Bài toán} & \textbf{Tên đầy đủ} & \textbf{Đầu vào} & \textbf{Đầu ra} \\
\midrule
ATE  & Aspect Term Extraction & Câu & Danh sách cụm từ khía cạnh \\
APC  & Aspect Polarity Classification & Câu + khía cạnh & Cảm xúc, nhóm khía cạnh \\
ASTE & Aspect Sentiment Triplet Extraction & Câu & (khía cạnh, ý kiến, cảm xúc) \\
ASQP & Aspect-Category-Opinion-Sentiment Quad & Câu & Bộ bốn ACOS \\
\bottomrule
\end{tabular}
\end{table}

\subsection{2.1.2.\; Bài toán trích xuất bộ ba (Joint Triplet Extraction)}
\label{subsec:triplet}

Khóa luận hình thức hóa đầu ra hệ thống là tập bộ ba:
\begin{equation}
\mathcal{T} = \{(a_i, c_i, s_i)\}_{i=1}^{m}
\label{eq:triplet}
\end{equation}
với $a_i$ là \textit{aspect term}, $c_i \in$ \{\texttt{AMENITY, BRANDING,
EXPERIENCE, FACILITY, LOYALTY, SERVICE}\} là \textit{aspect category},
và $s_i \in$ \{Positive, Negative, Neutral\}. Một bộ ba đúng khi
\textbf{cả ba thành phần} khớp đồng thời với nhãn gold.

% =====================================================================
\section{2.2.\; Kiến trúc Transformer và cơ chế Attention}
\label{sec:transformer}

Transformer \cite{vaswani2023attentionneed} thay thế cơ chế hồi quy
bằng \textbf{self-attention}, xử lý toàn chuỗi song song. Mỗi khối
encoder gồm: \textbf{Multi-Head Self-Attention (MHSA)} + residual +
LayerNorm, và \textbf{Feed-Forward Network (FFN)} + residual + LayerNorm.

\textbf{Self-Attention} chiếu chuỗi đầu vào $X \in \mathbb{R}^{n \times d}$
thành Query, Key, Value:
\begin{equation}
Q = XW^Q,\quad K = XW^K,\quad V = XW^V
\end{equation}
Đầu ra:
\begin{equation}
\text{Attention}(Q,K,V) = \text{softmax}\!\left(\frac{QK^\top}{\sqrt{d_k}}\right)V
\label{eq:attention}
\end{equation}

Chi phí tính toán tỉ lệ $O(n^2 \cdot d_k)$ -- đây là động lực cốt
lõi của Token Merging: giảm $n$ sẽ giảm chi phí bậc hai.

\textbf{Multi-Head Attention} dùng $h$ đầu song song, mỗi đầu học
một loại quan hệ khác nhau:
\begin{equation}
\text{MultiHead}(Q,K,V) = \text{Concat}(\text{head}_1,\ldots,\text{head}_h)W^O
\end{equation}

\textbf{FFN} hai lớp với hàm kích hoạt GELU:
\begin{equation}
\text{FFN}(x) = \text{GELU}(xW_1+b_1)W_2+b_2
\end{equation}

% =====================================================================
\section{2.3.\; Mô hình tiền huấn luyện: BERT và T5}
\label{sec:plm}

\subsection{2.3.1.\; BERT}
\label{subsec:bert}

BERT \cite{devlin2019bert} là encoder Transformer hai chiều, huấn luyện
trước bằng \textbf{Masked Language Modeling} (che 15\% token, dự đoán
lại từ ngữ cảnh hai chiều) và \textbf{Next Sentence Prediction}.

Tokenizer \textbf{WordPiece} thêm token đặc biệt: \texttt{[CLS]} (đầu
chuỗi, dùng cho phân loại) và \texttt{[SEP]} (phân tách). Đầu ra là
ma trận hidden states $H \in \mathbb{R}^{L \times 768}$ -- đầu vào
trực tiếp cho các module LCF và ToMe của khóa luận.

\noindent\textit{Vai trò trong repo}: \texttt{[CLS]} và \texttt{[SEP]}
được bảo vệ tuyệt đối (\texttt{protect\_cls=True, protect\_sep=True})
khỏi bị gộp trong \texttt{ToMeSequenceMerger}.

\subsection{2.3.2.\; T5}
\label{subsec:t5}

T5 \cite{raffel2020t5} (Text-to-Text Transfer Transformer) biểu diễn
mọi bài toán NLP dưới dạng \emph{sinh văn bản}. T5 theo kiến trúc
encoder--decoder, huấn luyện trước với mục tiêu \textit{span corruption}.

\noindent\textit{Vai trò trong repo}: \texttt{T5AspectExtractor}
(backbone \texttt{t5-base}) giải bài toán ATE bằng \textbf{sinh có điều
kiện} theo định dạng GAS -- câu đầu vào $\to$ danh sách khía cạnh
\texttt{"(term1); (term2)"} hoặc \texttt{"none"}. Đầu ra được chuẩn
hóa về vị trí trong câu gốc bằng đối sánh Levenshtein (mục~\ref{subsec:lev_ngram}).

\subsection{2.3.3.\; Giải mã Beam Search}
\label{subsec:beam_search}

\subsubsection{2.3.3.1.\; Vấn đề của greedy decoding}

Trong bài toán sinh chuỗi (sequence generation), \textbf{greedy decoding}
tại mỗi bước chọn token có xác suất cao nhất:
\begin{equation}
\hat{y}_t = \arg\max_{v \in \mathcal{V}} P(v \mid y_{<t}, x)
\label{eq:greedy}
\end{equation}
Cách này nhanh nhưng dễ rơi vào cực tiểu cục bộ: một lựa chọn tốt ở
bước $t$ có thể dẫn đến chuỗi kém ở các bước sau. Trong sinh danh sách
khía cạnh, greedy decoding có thể bỏ lỡ cụm từ ngắn hơn nhưng phổ biến
hơn do bị ``kéo'' về phía một nhánh sinh sớm.

\subsubsection{2.3.3.2.\; Thuật toán Beam Search}

\textbf{Beam Search} duy trì $k$ ứng viên (beam) song song tại mỗi bước
giải mã. Tại bước $t$, mỗi ứng viên $b_j$ được mở rộng thành $|\mathcal{V}|$
nhánh; trong tổng số $k \cdot |\mathcal{V}|$ ứng viên mới, chỉ giữ lại
$k$ ứng viên có log-xác suất tích lũy cao nhất:
\begin{equation}
\text{score}(y_{1:t}) = \sum_{\tau=1}^{t} \log P(y_\tau \mid y_{<\tau}, x)
\label{eq:beam_score}
\end{equation}
\begin{equation}
\mathcal{B}_t = \text{Top-}k\!\left\{
  (b, v) : b \in \mathcal{B}_{t-1},\; v \in \mathcal{V}
\right\}_{\text{score}}
\label{eq:beam_topk}
\end{equation}

Sau $T$ bước, ứng viên có score cao nhất trong tất cả beam là kết quả.
Với $k=1$, Beam Search suy biến về greedy decoding.

\textbf{Length penalty} thường được áp dụng để tránh thiên vị các chuỗi
ngắn:
\begin{equation}
\text{score\_norm}(y) = \frac{\text{score}(y)}{|y|^\alpha},\quad \alpha \in (0,1]
\label{eq:length_penalty}
\end{equation}

\noindent\textit{Vai trò trong repo}: \texttt{T5AspectExtractor} dùng
Beam Search với $k=4$ (\texttt{num\_beams=4}), \texttt{max\_length=64}.
Bốn beam cho phép mô hình khám phá đồng thời các cách ghép term khác
nhau (ví dụ \texttt{"(room)"} vs \texttt{"(the room)"}), giảm rủi ro
bỏ sót khía cạnh do greedy chọn nhầm token đầu tiên.

\subsection{2.3.4.\; Chuẩn hóa đầu ra bằng khoảng cách Levenshtein}
\label{subsec:levenshtein}

\subsubsection{2.3.4.1.\; Vấn đề lệch chuỗi sinh}

T5 sinh tự do không đảm bảo đầu ra là chuỗi con chính xác của câu gốc:
mô hình có thể sinh sai chính tả, đổi viết hoa/thường, thêm/bớt
whitespace, hoặc dùng dạng từ biến thể (ví dụ: \textit{``staffs''}
thay vì \textit{``staff''}). Nếu dùng trực tiếp chuỗi sinh làm aspect
term, bước APC sẽ không tìm được vị trí token tương ứng trong câu
$\Rightarrow$ \texttt{lcf\_vec} sai.

\subsubsection{2.3.4.2.\; Khoảng cách Levenshtein}
\label{subsec:levenshtein_dist}

\textbf{Khoảng cách chỉnh sửa Levenshtein} (edit distance) $\text{lev}(s, t)$
giữa hai chuỗi $s$ và $t$ là số thao tác \emph{tối thiểu} (chèn, xóa,
thay thế một ký tự) cần thiết để biến $s$ thành $t$:
\begin{equation}
\text{lev}(s,t) =
\begin{cases}
|s| & \text{nếu } |t|=0 \\
|t| & \text{nếu } |s|=0 \\
\text{lev}(s[1:], t[1:]) & \text{nếu } s[0]=t[0] \\
1 + \min \left\{
\begin{array}{l}
\text{lev}(s[1:], t) \\
\text{lev}(s, t[1:]) \\
\text{lev}(s[1:], t[1:])
\end{array}
\right. & \text{nếu } s[0]\neq t[0]
\end{cases}
\label{eq:levenshtein}
\end{equation}

Tính toán hiệu quả bằng lập trình động trong $O(|s|\cdot|t|)$.

\subsubsection{2.3.4.3.\; Chuẩn hóa n-gram bằng Levenshtein}
\label{subsec:lev_ngram}

Thay vì so sánh chuỗi sinh $\hat{a}$ với toàn bộ câu $x$, hệ thống
liệt kê tất cả \textbf{n-gram liên tiếp} (với $n$ bằng số token của
$\hat{a}$) trong câu gốc, rồi chọn n-gram có khoảng cách Levenshtein
nhỏ nhất:
\begin{equation}
a^* = \arg\min_{g \in \text{ngrams}(x,\,|\hat{a}|)} \text{lev}(\hat{a},\, g)
\label{eq:lev_match}
\end{equation}

Nếu $\text{lev}(\hat{a}, a^*) = 0$, cụm từ khớp chính xác. Nếu lớn hơn
một ngưỡng (ví dụ $>0.4 \cdot |\hat{a}|$), cụm từ bị loại bỏ (có thể
là ảo giác -- hallucination của T5).

\begin{table}[H]
\centering
\caption{Ví dụ chuẩn hóa Levenshtein trên đầu ra T5}
\label{tab:lev_example}
\small
\begin{tabular}{@{}p{3.5cm}p{3.5cm}cp{4.0cm}@{}}
\toprule
\textbf{Câu gốc} & \textbf{T5 sinh ra} & \textbf{lev} & \textbf{Sau chuẩn hóa} \\
\midrule
\textit{``The room was great''} & \texttt{room} & 0 & \texttt{room} \cmark \\
\textit{``Excellent staffs at front desk''} & \texttt{staffs} & 1 & \texttt{staffs} \cmark \\
\textit{``The wifi is slow''} & \texttt{wi-fi} & 2 & \texttt{wifi} \cmark \\
\textit{``Clean and spacious''} & \texttt{location} & 6 & loại bỏ \xmark \\
\bottomrule
\end{tabular}
\end{table}

\noindent\textit{Vai trò trong repo}: hàm \texttt{normalize\_aspect\_term()}
trong \texttt{src/inference.py} áp dụng đối sánh n-gram Levenshtein
trên từng cụm từ trong đầu ra beam search của \texttt{T5AspectExtractor}.
Bước này đảm bảo aspect term được trả về luôn là chuỗi con hợp lệ
xuất hiện trong câu gốc, cho phép \texttt{build\_lcf\_vec()} xác định
vị trí token chính xác qua \texttt{offset\_mapping}.

% =====================================================================
\section{2.4.\; Cơ chế Local Context Focus (LCF)}
\label{sec:lcf}

\subsection{2.4.1.\; Động lực}

Trong câu đa khía cạnh, từ cảm xúc của một khía cạnh thường xuất hiện
gần khía cạnh đó. Phân loại dùng toàn bộ biểu diễn câu dễ bị nhiễu
từ các khía cạnh khác. LCF \cite{zeng2019lcf, yang2021lcfatepc} tạo
thêm một \textbf{luồng cục bộ} tập trung vào vùng xung quanh khía cạnh
và kết hợp với luồng toàn cục.

\subsection{2.4.2.\; Khoảng cách ngữ nghĩa tương đối (SRD)}
\label{subsec:srd}

Vector chỉ thị nhị phân \texttt{lcf\_vec} $\in\{0,1\}^{L}$ đánh dấu
vị trí sub-word thuộc cụm khía cạnh. Từ đó xác định tâm và nửa độ rộng:
\begin{equation}
\text{center} = \frac{a_{start}+a_{end}}{2},\qquad
\text{half} = \left\lfloor\frac{a_{end}-a_{start}+1}{2}\right\rfloor
\end{equation}
SRD của token tại vị trí $i$:
\begin{equation}
SRD_i = \max\!\Big(0,\,\big|i-\text{center}\big|-\text{half}\Big)
\label{eq:srd}
\end{equation}
Token có $SRD_i \le \alpha$ ($\alpha=5$, \texttt{SRD\_THRESHOLD}) thuộc
\textbf{ngữ cảnh cục bộ}.

\subsection{2.4.3.\; Hai biến thể: CDM và CDW}

\textbf{Context Dynamic Mask (CDM)} -- mặt nạ nhị phân:
\begin{equation}
m_i = \begin{cases} 1 & SRD_i \le \alpha\\ 0 & SRD_i > \alpha \end{cases}
\qquad O^{CDM} = H \odot m
\label{eq:cdm}
\end{equation}

\textbf{Context Dynamic Weight (CDW)} -- trọng số giảm dần liên tục (mặc định):
\begin{equation}
c_i = \begin{cases} 1 & SRD_i \le \alpha\\[4pt]
\dfrac{L-(SRD_i-\alpha)}{L} & SRD_i > \alpha \end{cases}
\qquad O^{CDW} = H \odot c
\label{eq:cdw}
\end{equation}

CDM cắt cứng -- giúp tập trung hơn cho nhãn thiểu số.
CDW giảm mượt -- giữ thêm thông tin toàn cục, gradient lan truyền tốt hơn.
Trong repo, \textbf{CDW là mặc định} (\texttt{use\_cdm=False}).

\subsection{2.4.4.\; Kết hợp đặc trưng (Feature Fusion)}
\label{subsec:fusion}

Luồng cục bộ $H_{local} = H \odot c$ qua self-attention \texttt{bert\_SA};
luồng toàn cục $H$ giữ nguyên. Hai luồng ghép và chiếu tuyến tính:
\begin{equation}
H_{fused} = \texttt{Linear}_{2d\to d}\big([H_{local}\,;\,H]\big)
\label{eq:fusion}
\end{equation}
$H_{fused}$ tiếp tục qua self-attention \texttt{bert\_SA\_}, pooler
(\texttt{[CLS]}), rồi vào hai đầu phân loại.

% =====================================================================
\section{2.5.\; Token Merging (ToMe)}
\label{sec:tome}

\subsection{2.5.1.\; Ý tưởng và thuật toán gốc}

Token Merging \cite{bolya2023tome} gộp các token có biểu diễn \emph{tương
đồng nhau} (đo bằng cosine similarity) thành một token đại diện bằng
phép trung bình, giảm độ dài chuỗi mà không thêm tham số và không cần
huấn luyện lại. Thuật toán gốc (bipartite) chia token thành hai nhánh
xen kẽ $A$, $B$; mỗi $a \in A$ ghép với $b \in B$ có cosine lớn nhất:
\begin{equation}
\text{sim}(a,b)=\frac{\hat{x}_a\cdot\hat{x}_b}{\|\hat{x}_a\|\|\hat{x}_b\|}
\label{eq:cosine}
\end{equation}
Hai token ghép cặp được hợp nhất:
\begin{equation}
x_{merged} = \tfrac{1}{2}(x_a+x_b)
\label{eq:merge_avg}
\end{equation}

\subsection{2.5.2.\; Ba chiến lược gộp trong khóa luận}
\label{subsec:tome_strategies}

Module \texttt{ToMeSequenceMerger} (\texttt{token\_merging/tome\_1d.py})
hiện thực ba chiến lược qua tham số \texttt{merge\_strategy}:

\begin{description}
  \item[\texttt{bipartite}] (mặc định) Bipartite cosine matching theo
    ToMe CVPR 2023. Chia pool nội bộ thành $A$/$B$ xen kẽ; mỗi $a \in A$
    ghép với $b \in B$ tương đồng nhất (mỗi $b$ tối đa một lần).

  \item[\texttt{sequential\_local}] Quét trái$\to$phải; mỗi token so
    sánh cosine với hàng xóm trái/phải gần nhất, gập về phía tương đồng
    hơn. Ranh giới vùng khía cạnh (\texttt{mid\_sep}: $\text{lcf}_j{<}0.5$
    và $\text{lcf}_{j+1}{>}0.5$) được bảo vệ, không bị gộp.

  \item[\texttt{attention\_weighted}] Xếp hạng token theo attention tăng
    dần; lặp lại chọn token attention thấp nhất (chưa loại, không được
    bảo vệ) rồi gộp vào hàng xóm cosine cao nhất. Bảo vệ: \texttt{[CLS]},
    \texttt{[SEP]}, token khía cạnh ($\text{lcf}{>}0.5$), token attention
    top 25\%.
\end{description}

Cả ba chiến lược đều bảo vệ \texttt{[CLS]}, \texttt{[SEP]} và token
vùng khía cạnh (\texttt{protect\_aspect=True}). Sau khi gộp, chỉ thị
khía cạnh lan truyền theo phép $\max$:
$\texttt{lcf}[src] = \max(\texttt{lcf}[src],\,\texttt{lcf}[dst])$.

\subsection{2.5.3.\; Hai chế độ sau gộp: resize và compact}
\label{subsec:tome_modes}

\begin{description}
  \item[\texttt{resize=True}] Sau khi gộp xuống $L'$, nội suy tuyến tính
    (\texttt{F.interpolate, mode="linear"}) cả hidden states và
    \texttt{lcf\_vec} trở lại độ dài gốc $L$. Đầu ra luôn có shape
    $(B, L, d_h)$ -- các lớp sau giữ nguyên kiến trúc. \textbf{Mục đích:}
    cô lập và đo ảnh hưởng của thao tác gộp lên chất lượng biểu diễn,
    không đo tăng tốc.

  \item[\texttt{resize=False} (compact)] Giữ nguyên chuỗi rút gọn độ
    dài $L' < L$, đệm theo $L'_{\max}$ của batch. Các lớp self-attention
    phía sau nhận chuỗi ngắn hơn -- đây là cơ chế thực sự giảm chi phí
    tính toán. \textbf{Hạn chế:} mỗi câu/cấu hình có $L'$ khác nhau,
    cần đồng bộ hóa batch và cập nhật attention mask.

  \item[\texttt{pre\_bert}] (hướng phát triển) Gộp token ngay ở mức
    embedding, \emph{trước} 12 lớp Transformer backbone. Đây là đòn bẩy
    tốc độ lớn nhất vì rút ngắn chuỗi cho toàn bộ phần $O(n^2)$ của
    backbone. Mã hỗ trợ sẵn (\texttt{\_pre\_bert\_merge}) nhưng chưa
    đưa vào bộ cấu hình chính.
\end{description}

Toàn bộ thực nghiệm của khóa luận sử dụng \textbf{resize=True}; compact
và pre-BERT ToMe là hướng phát triển (Chương~\ref{chap:ketluan}).

% =====================================================================
\section{2.6.\; Huấn luyện đa nhiệm và hàm mất mát}
\label{sec:multitask_theory}

\texttt{FastLcfBertMultiTask} dự đoán đồng thời Sentiment (3 lớp) và
Aspect Category (6 lớp). Hàm mất mát tổng:
\begin{equation}
\mathcal{L} = w_s\,\ell_{\text{CE}}(f_s(h),y^{(s)})
  + w_c\,\mathbbm{1}[\neg\text{supp}]\cdot\ell_{\text{CE}}(f_c(h),y^{(c)})
\label{eq:mt_loss}
\end{equation}
Mẫu \textit{supplement} (\texttt{is\_supplement=True}) chỉ tham gia
hạng thứ nhất -- đầu phân loại category không học trên mẫu bổ trợ.
Cả hai thành phần dùng \textbf{trọng số lớp nghịch đảo tần suất}
$w_k = N/(C\cdot n_k)$ để giảm ảnh hưởng mất cân bằng.

% =====================================================================
\section{2.7.\; Độ đo đánh giá}
\label{sec:metrics_theory}

\begin{itemize}
  \item \textbf{ATE F1}: exact match cụm từ khía cạnh (Precision/Recall/F1
        ở mức micro trên cặp (câu, term)).
  \item \textbf{Micro F1 (bộ ba)}: TP/FP/FN tổng hợp toàn cục -- nhạy
        với nhãn phổ biến. Một bộ ba là TP khi $a=a^*\land c=c^*\land s=s^*$.
  \item \textbf{Macro F1}: trung bình F1 từng lớp ghép -- phản ánh nhãn
        thiểu số.
  \item \textbf{Oracle Joint Acc}: cho trước aspect term gold, đo
        \emph{trần} của bộ phân loại, tách biệt lỗi ATE.
  \item \textbf{Recall}: dùng ở Bước 5 -- tỉ lệ bộ ba gold được tìm thấy.
\end{itemize}

% =====================================================================
\section{2.8.\; Các hướng tiếp cận liên quan}
\label{sec:related_work}

\paragraph{APC với Transformer.} BERT \cite{devlin2019bert} đặt nền
tảng cho ABSA dạng SPC. FAST-LCF-BERT \cite{zeng2019lcf, li2020bertabsa}
bổ sung luồng ngữ cảnh cục bộ, đạt kết quả state-of-the-art. GAS
\cite{zhang2021gas} chuyển sang sinh bộ ba một bước bằng T5 -- module
\texttt{gas/} trong repo là biến thể liên quan.

\paragraph{Token Merging.} ToMe \cite{bolya2023tome} áp dụng cho Vision
Transformer; khóa luận chuyển thể cho NLP 1D với bổ sung bảo vệ token
ngữ nghĩa (khía cạnh, CLS) và đề xuất hai chiến lược mới
(\texttt{sequential\_local}, \texttt{attention\_weighted}).

\paragraph{Tách câu.} UOS (\texttt{uos/}) tách câu dài bằng LLM trước
inference, cùng hướng thu hẹp ngữ cảnh với LCF nhưng ở mức văn bản thô.

% =====================================================================
\section*{Tóm tắt chương}
\addcontentsline{toc}{section}{Tóm tắt chương}

Chương 2 trình bày cô đọng: (1) bài toán ABSA và cấu trúc bộ ba Joint
Triplet; (2) Transformer -- nền tảng tính toán và nguồn gốc chi phí
$O(n^2)$; (3) BERT (APC) và T5 (ATE) kèm kỹ thuật giải mã Beam Search
($k=4$) và chuẩn hóa Levenshtein đảm bảo aspect term hợp lệ; (4) LCF
với SRD, CDM/CDW và bước Fusion; (5) Token Merging với ba chiến lược
và ba chế độ sau gộp (resize/compact/pre-BERT); (6) hàm mất mát đa
nhiệm và các độ đo. Chương 3 trình bày chi tiết hiện thực và ví dụ
minh họa.
